\section{Occupation, exclusion and finite statistics}
\label{chap:ocupacion-estadistica}
\index{Pauli exclusion principle}\index{CAR}
\index{Fermi--Dirac}\index{Bose--Einstein}

The null and material realizations provide the geometric support of occupation. The fermionic sector is constructed on the canonical anticommutation relations (CAR), and the spin--statistics theorem further requires the declared operator representation. This section fixes those premises before deriving idempotence of occupation, the microcanonical cardinalities, and the Bose--Einstein and Fermi--Dirac distributions with their common classical limit.

\subsection{Idempotency of Occupation in the Anticommutation Algebra}

Let \(a_i,a_i^\dagger\) be operators satisfying the canonical
anticommutation relations (CAR)
\[
 \{a_i,a_j^\dagger\}=\delta_{ij},
 \qquad \{a_i,a_j\}=0,
\]
and let \(n_i=a_i^\dagger a_i\).

\begin{theorem}[Pauli Exclusion Principle]
\label{thm:principio-exclusion-pauli}
One has \(n_i^2=n_i\), and therefore
\(\Spec(n_i)\subseteq\{0,1\}\).
\end{theorem}
\begin{proof}
\[
 n_i^2=a_i^\dagger a_i a_i^\dagger a_i
 =a_i^\dagger(1-a_i^\dagger a_i)a_i
 =a_i^\dagger a_i,
\]
because \(a_i^2=(a_i^\dagger)^2=0\). Every eigenvalue of an idempotent
satisfies \(\lambda^2=\lambda\) \cite{Pauli1925,Dirac1926}.
\end{proof}

The nilpotence \((a_i^\dagger)^2=0\) precludes double fermionic occupation of the
same mode and, upon maximizing entropy under the energy and number
constraints, leads to the Fermi--Dirac occupation law presented below.
The symmetric completion and the Bose--Einstein law constitute a distinct
stratum and are presented separately.

The CAR relations are an algebraic hypothesis of the theorem. A monodromy
with sign change under a rotation of \(2\pi\) is compatible with a
fermionic representation, but does not by itself constitute a derivation of
CAR or of the spin--statistics theorem.
